
\documentclass[11pt]{article}
\usepackage[margin=1in]{geometry}
\usepackage[T1]{fontenc}
\usepackage[utf8]{inputenc}
\usepackage{lmodern}
\usepackage{setspace}
\usepackage{enumitem}
\usepackage{hyperref}
\usepackage{xcolor}
\hypersetup{colorlinks=true, linkcolor=blue, urlcolor=blue, citecolor=blue}
\setstretch{1.08}
\setlength{\parskip}{0.6em}
\setlength{\parindent}{0pt}

\begin{document}

\begin{center}
    {\Large \textbf{Memo for Revising the Paper for RAST Submission}}\\[0.5em]
    {\normalsize Prepared for the purpose of sharpening fit, focus, and execution for a submission to the Review of Accounting Studies conference on investor information needs}
\end{center}

\vspace{0.75em}

\textbf{Bottom line.} The paper can fit RAST well, but only if it is presented as a focused \emph{investor information-processing paper in an auditing setting}, not as a paper that tries to make five separate contributions at once. The revision objective should be to make the paper feel \emph{tighter}, not broader. The strongest version of the paper asks one central question: \emph{how do investors interpret an ambiguous, mandatory audit-related disclosure, and why does that interpretation vary systematically with polarization?}

\section*{1. Core positioning for RAST}

The paper should be framed around a single main contribution:
\begin{quote}
Auditor-change disclosures transmit information to investors, but investors do not interpret that information uniformly; polarization magnifies heterogeneity in interpretation and therefore magnifies market reactions.
\end{quote}

That framing works for RAST because it connects directly to investor information use, heterogeneous interpretation of disclosure, and the design limits of disclosure-based regulation. It also preserves the auditing identity of the paper. The revised draft should avoid sounding like a general paper on politics and markets with auditing used only as a convenient setting.

\textbf{Recommended one-sentence summary for the introduction.}

\begin{quote}
We examine how investors interpret mandatory auditor-change disclosures and show that disagreement in that interpretation is stronger in more polarized environments, leading to stronger capital-market reactions to the same audit-related signal.
\end{quote}

\section*{2. What to emphasize}

The following points should be elevated in the introduction, abstract, and conclusion.

\textbf{A. Investor interpretation of disclosure.} The paper is not merely about whether auditor changes matter; it is about how investors process a disclosure that is mandatory, salient, and often ambiguous in meaning.

\textbf{B. Heterogeneous users of the same information.} The central insight for RAST is that the same disclosure can generate different inferences across investors or investor environments. That is the bridge from auditing to the conference theme.

\textbf{C. Limits of disclosure alone.} The paper can make a restrained but useful regulatory point: more disclosure does not guarantee uniform understanding when interpretation depends on prior beliefs or social context.

\textbf{D. Auditor-change disclosures as part of the information environment.} The draft should say clearly that auditor-change disclosures matter because they inform investors about the credibility of the firm’s reporting environment and the reliability of reported accounting information going forward.

\section*{3. What not to do}

To keep the paper crisp and avoid the appearance of overreach, the following moves should be avoided.

\textbf{A. Do not turn the paper into a broad political-finance paper.} Polarization should remain a mechanism that affects interpretation of audit-related information, not the paper’s primary subject.

\textbf{B. Do not claim too many contributions.} The paper does not need to be simultaneously sold as a leading paper on auditing, investor attention, political economy, standard-setting, and social disagreement. That will make the draft feel scattered.

\textbf{C. Do not overstate the standard-setter angle.} A short implications paragraph is helpful. A sweeping claim about how standard setters should redesign disclosure regimes without direct evidence would weaken the paper.

\textbf{D. Do not imply evidence on investor groups that the paper does not directly observe.} It is acceptable to discuss heterogeneous interpretation; it is not acceptable to imply direct identification of which exact investor types drive the effect unless the empirics support that claim.

\section*{4. Concrete revision plan}

\textbf{Step 1: Rewrite the abstract.} The abstract should do three things in sequence:
\begin{enumerate}[leftmargin=1.5em]
    \item Introduce the setting as a mandatory audit-related disclosure with ambiguous information content.
    \item State the core empirical finding that market reactions are stronger in more polarized environments.
    \item End with the implication that disclosure effectiveness depends not only on content, but also on how investors interpret that content.
\end{enumerate}

\textbf{Step 2: Rewrite the first three pages of the introduction.} The introduction should be streamlined into the following flow:
\begin{enumerate}[leftmargin=1.5em]
    \item Why auditor-change disclosures matter for investors.
    \item Why such disclosures are especially open to heterogeneous interpretation.
    \item Why polarization is a plausible and economically meaningful determinant of that heterogeneity.
    \item What the paper finds.
    \item Why the findings matter for the literature on investor use of accounting and audit-related information.
\end{enumerate}

The contribution paragraph should be shortened and disciplined. Ideally it should contain \emph{three} contributions, not more:
\begin{enumerate}[leftmargin=1.5em]
    \item evidence on capital-market reactions to auditor-change disclosures,
    \item evidence that polarization amplifies heterogeneity in those reactions,
    \item implications for the effectiveness of disclosure-based information transmission.
\end{enumerate}

\textbf{Step 3: Tighten the theory-to-empirics mapping.} The theory section should make only the claims needed to support the empirical tests. If the model most naturally predicts stronger disagreement and trading volume, then that should be stated plainly. If the link to return magnitude depends on additional assumptions, those assumptions should be stated. The paper should not sound as if the formal model cleanly proves more than it actually does.

\textbf{Step 4: Strengthen the empirical motivation for the polarization measure.} The introduction and variable-design discussion should directly acknowledge the likely reader concern that headquarters-state polarization is an imperfect proxy for the investor environment. The draft should then explain, with discipline, why the proxy is still informative in this setting and what interpretation the paper is and is not claiming.

\textbf{Step 5: Add a short concluding paragraph tailored to RAST.} The conclusion should not merely restate the main result. It should explain that the evidence suggests that the informational effects of disclosure depend on the social environment in which investors interpret the signal. That makes the paper conference-relevant without requiring the draft to become a paper about standard setting per se.

\section*{5. Suggested language for the introduction}

The following sentence-level shifts would help.

Instead of:
\begin{quote}
Our paper studies how political polarization affects market reactions to auditor changes.
\end{quote}

Prefer:
\begin{quote}
Our paper studies how investor interpretation of mandatory auditor-change disclosures varies with polarization, and how that variation shapes the capital-market consequences of audit-related information.
\end{quote}

Instead of:
\begin{quote}
We contribute to the literature on political polarization in financial markets.
\end{quote}

Prefer:
\begin{quote}
We contribute to research on how investors use accounting and audit-related disclosures by showing that the same mandatory disclosure can generate systematically different inferences across more and less polarized environments.
\end{quote}

Instead of:
\begin{quote}
These findings have important implications for regulators.
\end{quote}

Prefer:
\begin{quote}
These findings suggest limits to disclosure-based regulation when users bring different priors to the interpretation of the same audit-related signal.
\end{quote}

\section*{6. Recommended tone for the paper}

The revised draft should sound \emph{confident but narrow}. That means:
\begin{itemize}[leftmargin=1.5em]
    \item do not oversell novelty where the evidence is still developing;
    \item do not claim universal conclusions beyond the auditor-change setting;
    \item do not over-interpret the mechanism beyond what the model and tests can sustain;
    \item do make the setting feel important, because auditor changes affect perceptions of reporting credibility and governance.
\end{itemize}

A good rule is that every major section should answer the same question: \emph{how does polarization change the way investors interpret an audit-related disclosure?} If a paragraph does not help answer that question, it should probably be cut or shortened.

\section*{7. Priority ordering}

If time is limited before submission, the revision priorities should be:

\begin{enumerate}[leftmargin=1.5em]
    \item rewrite the abstract and introduction around investor interpretation of audit-related disclosure;
    \item align the theory claims tightly with the empirics;
    \item discipline the contribution claims to three concise points;
    \item improve the discussion of the polarization proxy and its interpretation;
    \item add a restrained RAST-oriented implications paragraph in the conclusion.
\end{enumerate}

\section*{8. Final recommendation}

The right strategy is \textbf{not} to add more topics. The right strategy is to make the paper more coherent. For RAST, that means presenting the study as a focused paper on \emph{how investors use and interpret audit-related disclosures under heterogeneous social conditions}. If that framing is executed cleanly, the paper should read as more relevant to the conference while still remaining disciplined enough for eventual journal submission.

\end{document}
